% ============================================================
% §6  Experimental Setup  (§6.1 datasets + degradation protocol + metrics /
%      §6.2 enhancement baselines)
% Draft -- Writer, session (W5 reorder). Reframed from the old main.tex
% "§ The ITB Benchmark" (project/meeting/ICLR2027/main.tex L250-269).
%
% Drift contract (STORY S6 + ACCEPTANCE C1.4 + R-rules):
%   * The benchmark is DEMOTED from a contribution to experimental setup.
%     We DELETE the old "25-lever / 9-baseline-as-contribution" framing
%     (old L263-269): the stratified protocol is described here as the
%     evaluation substrate, not sold as a deliverable.
%   * §6.2 lists exactly SIX SOTA enhancement baselines (ACCEPTANCE C1.4):
%     Restormer / NAFNet / MIRNet-v2 / SwinIR / Uformer-B / Real-ESRGAN.
%     DIFFUSION IS FORBIDDEN (medical-safety red line): no DiffBIR / SD-Turbo.
%   * Degradation protocol numbers are taken from data/degrade.py (medium tier =
%     the "moderate" protocol used by E1-E12). Numbers not directly read from the
%     pipeline are left as \todo{verifier} rather than guessed.
%   * R-rules: every reported metric carries a bootstrap 2000-sample 95% CI (R6);
%     no SOTA claims here (R5); q-bar is interchangeable (R7); BMVC cited not
%     reused (R10); anonymise before submission (R4 -- ANON-* macros).
% ============================================================

\section{Experimental Setup}
\label{sec:setup}

\subsection{Datasets, degradation protocol, and metrics}
\label{subsec:setup-data}

\paragraph{Datasets.}
We evaluate on dermoscopic skin-lesion imagery drawn from the public ISIC
collection, using a fixed train/validation/test split; the held-out test split
contains $n=3627$ images with $117$ melanoma-positive cases, and the
diagnostic backbone is an EfficientNet-B3 classifier (denoted B3) frozen across
all evaluations so that any change in diagnostic AUC is attributable to the
input, not to retraining.\todo{verifier: confirm dataset source pool, split
sizes, and B3 backbone identity against DATA\_INVENTORY / training config}
All imagery is public; we conduct no reader study and no offline collection.

\paragraph{Degradation protocol.}
To study quality-conditioned behaviour we apply controlled, measured degradations
rather than relying on naturally occurring low-quality images, so that quality is
a known variable. Five physically parameterised dimensions are modelled --- blur
(Gaussian $\sigma$), brightness (multiplicative factor), completeness (retained
crop ratio), colour shift (per-channel offset), and contrast (scaling
$\alpha$) --- mirroring the failure modes of consumer-device capture. Unless
otherwise stated, the main paired-evaluation protocol applies the \emph{moderate}
degradation tier (blur $\sigma=1.5$; brightness factor sampled in $[0.65, 0.85]$;
crop ratio in $[0.75, 0.89]$; colour-shift offset $0.12$; JPEG quality in
$[50, 74]$), composited at $256\times256$, followed by enhancement at $256\times256$,
a centre crop to $224\times224$, and the B3 classifier.\todo{verifier: confirm
moderate-tier composition (activation probabilities and JPEG range) against
data/degrade.py; confirm whether contrast/$\alpha$ is part of the moderate
paired protocol or only the per-dimension severity sweep} For the per-dimension
$\times$ per-severity decision surface of \S\ref{sec:decision-surface}, each
dimension is instead swept across five severity levels along its physical scalar
(severity grid in Table~\ref{tab:c0-severity}).

\paragraph{Metrics.}
Restoration fidelity is reported as per-image PSNR --- computed as the mean over
images of $10\log_{10}(1/\mathrm{MSE}_i)$, the per-image convention we use
throughout --- together with SSIM. Diagnostic behaviour is reported as the
diagnostic AUC of the frozen B3 classifier, the change in AUC induced by
enhancement ($\Delta\mathrm{AUC}$), the classification-agreement rate between
enhanced and reference predictions, and the feature-level KL divergence that the
DP-Loss controls. Every PSNR, AUC, $\Delta\mathrm{AUC}$, and agreement figure is
accompanied by a bootstrap $2000$-sample $95\%$ confidence interval, and paired
comparisons additionally report a McNemar test. Inference latency is reported in
milliseconds per image. The companion calibration study reports a complementary
per-dimension expected-calibration-error and quality-stratified view; we cite it
\citep{ANON-BMVC} but do not reuse its tables here.

\subsection{Enhancement baselines}
\label{subsec:setup-baselines}

We compare our diagnosis-preserving enhancement against six representative,
publicly available restoration backbones spanning convolutional and transformer
designs: \textbf{Restormer}~\citep{zamir2022restormer},
\textbf{NAFNet}~\citep{chen2022simple},
\textbf{MIRNet-v2}~\citep{zamir2022mirnetv2},
\textbf{SwinIR}~\citep{liang2021swinir},
\textbf{Uformer-B}~\citep{wang2022uformer}, and
\textbf{Real-ESRGAN}~\citep{wang2021realesrgan}. These cover the strongest
deterministic restoration families in current use. We deliberately exclude
diffusion- and other generative restoration models: on dermoscopy such models
can hallucinate diagnostically salient structure (invented pigment networks or
vascular patterns), which is a medical-safety hazard rather than a desirable
prior; we analyse this failure mode explicitly in \S\ref{sec:discussion} rather
than admit it as a baseline. Each baseline is run under the identical degradation
protocol and the same frozen B3 classifier described above, so that the
comparison isolates the effect of the enhancement model on diagnostic
preservation; results appear in \S\ref{subsec:setup-baselines}
(Table~\ref{tab:e10}).\todo{verifier: confirm each baseline's checkpoint /
training configuration source (official release vs. re-trained) against the
experiment log for E10, job 1448952}
